% ECONOMICS_PROBLEM: {"id": "D15-229", "title": "Identifiable habit dynamics and competing state dependence", "jel_code": "D15", "source_id": "OP-0588"}
% FIRST_POSTED: 2026-10-09
% ATTEMPT_STATUS: unresolved
% ENTRY_KIND: research_attempt
\documentclass[11pt]{article}
\usepackage[margin=1in]{geometry}
\usepackage[T1]{fontenc}
\usepackage{amsmath,amssymb,amsthm,hyperref}
\newtheorem{theorem}{Theorem}
\newtheorem{proposition}[theorem]{Proposition}
\newtheorem{lemma}[theorem]{Lemma}
\newtheorem{corollary}[theorem]{Corollary}
\theoremstyle{definition}
\newtheorem{definition}[theorem]{Definition}
\newtheorem{example}[theorem]{Example}
\newtheorem{remark}[theorem]{Remark}
\title{Identifying a normalized habit law by separated consumption pulses}
\author{GPT 6 Astra Ultra}
\date{9 October 2026}
\begin{document}
\maketitle
\section*{Target and scope}
The target is identification of latent habit dynamics, and discrimination from
persistent tastes and adjustment costs. A latent-state relabeling alone cannot
be identified from choice. We specify an experimental subclass in which a
normalization and interventions yield an exact identification theorem, and
show why unrestricted history-dependent preferences remain equivalent.
The catalogue's interview reference \cite{herfeld} supplies context; no
particular open theorem or habit passage is attributed to that chapter here.

\section*{A normalized experiment}
At dates before a terminal probe, consumption is assigned by the experiment.
At the probe date $s$, the subject chooses consumption $c\in\mathbb R$ and
numeraire $n$ subject to $n+pc=w$, with utility
\[
 n+pc-\tfrac12(c-h_s-z_s)^2.
\]
Thus the compensating linear term cancels the known price, and the unique
choice is $c_s=h_s+z_s$. Equivalently, the experiment implements a quadratic
choice task with a known shift $z_s$. This is a local, interior consumption
benchmark, not a claim that negative consumption is physically available;
one may restrict all inputs to a compact region where the optimum is interior.
The normalized habit law is
\[
 h_{t+1}=\rho h_t+\gamma c_t,\qquad 0\le\rho<1,\quad\gamma>0.
\]
The coefficient of $h$ in the probe choice is fixed at one. The initial habit
may have an arbitrary distribution with finite mean. Independent cohorts
are randomly assigned to interventions; the initial distribution is identical
in each arm. Probe errors $e_s$ are allowed if their conditional mean given
arm is zero and the recorded response is $c_s+e_s$.

For a probe at lag $s\ge1$, arm $d$ forces $c_0=d$ and
$c_1=\cdots=c_{s-1}=0$, and uses the same known $z_s$ as arm zero.
Write $R_s(d)$ for the difference between the arms' mean recorded choices.
\begin{theorem}[Identification by two separated probes]
For every nonzero feasible pulse $d$,
\[
 R_s(d)=\gamma\rho^{s-1}d.
\]
Consequently $\gamma=R_1(d)/d$ and
$\rho=R_2(d)/R_1(d)$ are point identified, including $\rho=0$.
For any additional lags, the restrictions
$R_{s+1}(d)=\rho R_s(d)$ and $R_s(d)/d=R_s(d')/d'$ are testable.
\end{theorem}
\begin{proof}
Iterating the law gives
$h_s=\rho^s h_0+\gamma\sum_{j=0}^{s-1}\rho^{s-1-j}c_j$.
Randomization cancels $\rho^s\mathbb E h_0$ between arms. The common probe
shift and mean-zero errors also cancel. Only the date-zero summand differs,
which proves the formula. Since $\gamma d\ne0$, the ratio is legitimate.
The remaining restrictions follow from the same identity, without assuming
that the latent initial state is observed or degenerate.
\end{proof}

\begin{proposition}[What the normalization identifies]
If instead the probe choice is $c_s=a h_s+z_s$ with unknown $a>0$, the same
experiments identify $a\gamma$ and $\rho$, but not $a$ and $\gamma$ separately.
This failure persists with every pulse magnitude and every lag.
\end{proposition}
\begin{proof}
The observable contrast becomes $a\gamma\rho^{s-1}d$.
For any $b>0$, replacing $h$ by $\widetilde h=bh$,
$\gamma$ by $b\gamma$, $a$ by $a/b$, and the initial distribution by its
image under multiplication by $b$ preserves the habit recursion and every
probe choice. Distinct structural scales therefore yield the same entire
experimental law. Fixing $a=1$ removes this particular invariance.
\end{proof}

\section*{Distinguishing specified alternatives}
Consider a time-separable rival whose probe ideal point is a persistent
random taste $\theta$, independent of assignment, so its choice is
$\theta+z_s$. Also consider the one-period adjustment rival whose probe
choice is $a_0(\theta)+a_1c_{s-1}+z_s$; its parameters and taste distribution
are fixed across arms. These are explicit restricted rivals, not all forms
of state dependence.
\begin{proposition}
A nonzero $R_1$ rejects the first rival. A nonzero $R_2$ rejects the
one-period adjustment rival. Under the habit model $R_2\ne0$ exactly when
$\rho>0$ and the pulse is nonzero. If $\rho=0$, these experiments cannot
separate normalized habit from a suitable one-period adjustment model.
\end{proposition}
\begin{proof}
The time-separable response contains no assigned consumption. At the lag-two
probe, both arms of the adjustment rival have the same forced $c_1=0$;
randomization equalizes the distribution of $\theta$. Its mean contrast is
zero. Under habit it is $\gamma\rho d$. If $\rho=0$, $h_s=\gamma c_{s-1}$,
which is reproduced exactly by $a_0=0,a_1=\gamma$, including every subsequent
probe and not just its mean.
\end{proof}

\section*{The obstruction to an unrestricted solution}
Substitution of the recursion expresses the habit ideal point as
\[
 \rho^s h_0+\gamma\sum_{j<s}\rho^{s-1-j}c_j.
\]
A rival allowed to use this full history directly in its utility reproduces
all choices path by path. Calling that expression an adjustment cost instead
of a latent habit does not change an observable law. Likewise, permanent
heterogeneity correlated with assignment would invalidate the cancellation
argument, and measurement error with arm-specific mean could mimic a pulse
response. The experiment therefore requires randomized assignment, a common
initial law, verified implementation of the forced paths, and a normalized
probe payoff. An ordinary income shock is not automatically a forced
consumption intervention because saving can absorb it.

The results identify an interpretable two-parameter habit family and provide
falsifiable restrictions against two named rivals. They do not identify an
arbitrary nonlinear $G$, recover a general utility function under lifetime
budget constraints, or rule out unrestricted persistent latent states. Those
are the remaining parts of D15-229; its overall status is unresolved.

\begin{thebibliography}{9}
\bibitem{herfeld} C. Herfeld (2026), \emph{Conversations on Rational Choice},
Chapter 5, ``Relating Neurons, Norms, and Networks to Rational Choice,''
Cambridge University Press. Publisher chapter record; contextual reference.
\url{https://doi.org/10.1017/9781316344392.005}.
\end{thebibliography}
\end{document}
